\chapter{Background Literature}
\label{ch:background}

\section{Problem Domain: Urban Transportation and Ride-Sharing}
\label{sec:domain}

Urban transportation plays an important role in providing access to employment, education, services and other daily activities. In many urban areas, commuters depend on a combination of private vehicles and public or informal transport services to reach their destinations. The availability, coverage and flexibility of these transport options can influence how easily people are able to travel between different locations. \citet{cervero2007} explain that informal transport systems can provide important mobility services, particularly where formal public transport does not adequately meet travel demand. However, informal transport can also experience challenges associated with coordination, service coverage and the organisation of transport services.

These issues are relevant to the transportation environment in Botswana, where commuters may rely on different forms of transport depending on their location and destination. Local reporting has, for example, documented concerns from combi operators in Francistown regarding the extension of routes to better serve areas such as the industrial area \citep{kebareng2013}. Although this example does not represent the entire transportation system in Botswana, it demonstrates how route configuration and access to particular destinations can affect local transport operations. The need to connect people with destinations therefore extends beyond simply having vehicles available; the routes and travel patterns of both transport providers and passengers are also important.

Ride-sharing provides an alternative approach to transportation in which individuals travelling in similar directions can share a vehicle. Unlike conventional public transport, where vehicles generally follow predetermined routes, ride-sharing can use the journeys of individual vehicle owners to create flexible transportation opportunities. \citet{agatz2012} describe dynamic ride-sharing as a process in which drivers and passengers with similar itineraries and schedules can be brought together, allowing available vehicle capacity to be used more effectively. This makes ride-sharing particularly relevant to situations where private vehicles may have unused seating capacity.

The potential benefits of ride-sharing extend beyond simply providing another method of transportation. \citet{furuhata2013} identify reduced travel costs, improved utilisation of vehicle capacity and possible reductions in traffic and environmental impacts among the potential benefits associated with ridesharing. However, these benefits depend on the ability of a system to identify people whose journeys are sufficiently compatible. A driver travelling in one direction cannot necessarily accommodate a passenger travelling in another direction without introducing additional distance or time into the journey.

Consequently, an effective ride-sharing system must consider more than the availability of a vehicle. It must consider factors such as the driver's origin and destination, the passenger's origin and destination, departure times, available seats and the extent to which the two journeys can be combined. \citet{stiglic2016} demonstrate that driver and passenger flexibility can have an important influence on the number of successful ridesharing matches. This indicates that the effectiveness of a ride-sharing system is strongly connected to the way in which potential journeys are identified and matched.

The proposed project therefore focuses on a peer-to-peer ride-sharing model for urban commuting in Botswana. The intended model allows a private vehicle owner who is already travelling towards a destination to offer available seats to passengers travelling along a compatible route. Instead of functioning primarily as an on-demand transport service in which a driver is expected to travel specifically to collect a passenger, the proposed system is intended to make use of journeys that are already taking place. The central problem addressed by the project is therefore the difficulty of efficiently connecting available vehicle capacity with passengers whose journeys are sufficiently compatible to be shared.

This creates an opportunity for a digital platform to assist with the coordination of these journeys. By allowing users to provide their locations, destinations, travel times and seating requirements, the application can use this information to identify potentially suitable rides. The following section examines existing ride-sharing and ride-hailing platforms to determine which approaches are already being used and which aspects can inform the design of the proposed system.

\section{Existing Ride-Sharing and Ride-Hailing Systems}
\label{sec:existing}

The development of digital transportation platforms has demonstrated that mobile applications can be used to coordinate passengers, drivers, locations and journeys. Existing systems provide useful examples of how users can request transportation, identify available drivers, communicate journey information and manage trips. However, these systems do not all use the same transportation model. Understanding these differences is important because the proposed application is not intended to simply reproduce an existing ride-hailing platform.

Uber provides an example of a conventional on-demand ride-hailing model. A passenger enters a pickup location and destination, submits a ride request, and is matched with a nearby driver who accepts the request. The application can then provide information about the driver's arrival, the trip and navigation, while users can rate one another after the journey \citep{uber}. These features demonstrate the usefulness of combining user accounts, geolocation, ride requests, trip information and status updates within one mobile platform. The proposed system can therefore take inspiration from Uber's digital interaction between passengers and drivers, particularly its use of location information and a structured ride-request process.

However, the transportation model is different from the one proposed in this project. In a conventional ride-hailing system, the driver accepts a passenger's request and travels to the passenger's pickup location. The proposed system instead focuses on a driver who already has a journey planned and has available seats. Therefore, the central question in the proposed system is not simply whether a driver is available near a passenger, but whether an existing driver's journey is sufficiently compatible with the passenger's requested journey.

Bolt provides another example of an app-based ride-hailing platform. Bolt describes its service as connecting passengers with drivers and allows users to enter journey information, receive ride options and request transportation. The platform also supports scheduled rides in advance and provides estimated prices based on factors such as journey distance and selected ride type \citep{bolt}. From a system-design perspective, Bolt demonstrates how a transportation application can combine location-based requests, driver availability and journey information within a mobile interface. These concepts are relevant to the proposed application's user interface and ride-request workflow.

Nevertheless, Bolt remains primarily a ride-hailing service in which a passenger requests transportation from an available driver. The proposed application instead aims to coordinate shared journeys between private vehicle owners and passengers. This distinction is important because the matching mechanism must prioritise route compatibility and available seating rather than simply finding an available nearby vehicle.

inDrive demonstrates a somewhat different approach to the interaction between passengers and drivers. Users can specify a pickup location, destination and proposed price, after which drivers can respond to the request. Drivers can accept suitable requests, decline them or make counter-offers, while passengers can consider different driver responses before selecting a driver \citep{indrive}. This provides a useful example of allowing participants to have greater involvement in the process of arranging a journey rather than relying entirely on automatic assignment.

For the proposed system, the relevant concept is the interaction between a ride request and potential drivers. The proposed application can similarly use a request-and-response process in which a passenger requests a compatible ride and a driver can accept or decline the request. However, the proposed system does not intend to reproduce inDrive's price-negotiation model. Instead, the main consideration is whether the passenger's journey can be accommodated within the driver's existing route and available seating.

BlaBlaCar is particularly relevant because its model is closer to the peer-to-peer ride-sharing approach proposed in this project. BlaBlaCar connects drivers travelling towards a destination with passengers travelling in the same direction so that they can share the journey and associated costs. Drivers can publish journeys by specifying information such as departure and arrival locations, dates, times and available seats, while passengers can search for suitable journeys using criteria such as origin, destination and travel date \citep{blablacar}.

The BlaBlaCar model therefore provides a strong conceptual foundation for the proposed application. In particular, the concepts of available seats, shared journeys, compatible destinations and cost contribution are directly relevant to the project. BlaBlaCar has also developed approaches for shorter commuting journeys in which drivers and passengers can be matched according to similar itineraries and schedules, demonstrating the relevance of route and time compatibility to regular commuting \citep{blablacar}.

However, the proposed system is not intended to reproduce BlaBlaCar as a complete platform. The project focuses specifically on urban commuting in Botswana and is designed around a simplified peer-to-peer model suitable for the scope of the prototype. The system therefore takes inspiration from BlaBlaCar's core concept of connecting drivers with available seats to passengers travelling in a similar direction, while adapting the concept to the project's intended local commuting context.

\subsection*{Comparison of Existing Systems}

\begin{longtable}{@{}>{\raggedright\arraybackslash}p{0.14\textwidth}
                   >{\raggedright\arraybackslash}p{0.36\textwidth}
                   >{\raggedright\arraybackslash}p{0.42\textwidth}@{}}
\caption{Comparison of existing ride-sharing and ride-hailing systems}\label{tab:comparison}\\
\toprule
\textbf{Existing System} & \textbf{Relevant features and approach} & \textbf{Relevance to the proposed system} \\
\midrule
\endfirsthead
\toprule
\textbf{Existing System} & \textbf{Relevant features and approach} & \textbf{Relevance to the proposed system} \\
\midrule
\endhead
\bottomrule
\endfoot
Uber & On-demand ride requests, pickup and destination selection, driver matching, trip information, navigation and post-trip ratings. & Informs the use of user accounts, geolocation, ride requests, trip status and interaction between drivers and passengers. The proposed system differs by focusing on drivers who are already travelling and have available seats. \\
\addlinespace
Bolt & Mobile ride-hailing, location-based driver connection, journey information, ride options and scheduled trips. & Informs the request-and-response interaction between passengers and drivers. The proposed system instead prioritises route compatibility and available seats rather than price negotiation. \\
\addlinespace
inDrive & Passengers provide journey details and proposed prices, while drivers can respond, accept, decline or make counteroffers. & Informs the request-and-response interaction between passengers and drivers. The proposed system instead prioritises route compatibility and available seats rather than price negotiation. \\
\addlinespace
BlaBlaCar & Drivers publish journeys with available seats and passengers search for compatible journeys based on origin, destination, time and available seats. Costs are shared between participants. & Provides the closest conceptual model for the proposed system. The project adapts the available-seat and shared-journey approach for an urban commuting application focused on Botswana. \\
\end{longtable}

The comparison demonstrates that existing platforms provide several useful components for the proposed system, but they approach transportation in different ways. Uber and Bolt demonstrate the effectiveness of location-based on-demand ride requests, while inDrive demonstrates participant involvement in arranging rides. BlaBlaCar provides the closest example of the proposed peer-to-peer carpooling model because it connects drivers with available seats to passengers travelling in the same direction.

At the same time, the comparison identifies an important design opportunity. The proposed application is intended to combine selected concepts from these systems without becoming a conventional ride-hailing platform. Its main focus is on existing journeys, available vehicle capacity and route compatibility. A driver is therefore not primarily treated as a transport provider waiting for a passenger request; instead, the driver can offer seats on a journey they already intend to make. This distinction gives the system a different operational focus from conventional ride-hailing.

The proposed system is consequently intended to provide a localised alternative centred on peer-to-peer commuting and cost sharing. Its value proposition is not based on assuming that users will automatically prefer it or that it will always be cheaper than other transportation options. Rather, the system is designed to provide an alternative way of coordinating journeys by allowing available vehicle capacity to be shared between people travelling along compatible routes. Whether users consider this approach useful or preferable would need to be evaluated through system testing and user feedback.

The comparison also shows that a successful implementation cannot rely only on identifying nearby users. A passenger may be geographically close to a driver while travelling towards a completely different destination. Conversely, users may begin from different locations while travelling along a substantial portion of the same route. Therefore, the proposed application requires a matching mechanism capable of considering the relationship between complete journeys. The literature on ride-matching provides the basis for determining which factors should be considered, as discussed in the following section.

\section{Ride-Matching and Compatibility}
\label{sec:matching}

Ride-matching is the process of identifying drivers and passengers whose journeys are sufficiently compatible to be combined into a shared trip. It is a central component of a dynamic ride-sharing system because the usefulness of such a system depends on its ability to identify suitable journeys rather than simply connecting an available driver with any passenger. \citet{agatz2012} describe dynamic ride-sharing as a process of bringing together individuals with similar itineraries and schedules, while considering the available capacity of participating vehicles.

A basic matching process can consider the origin and destination of both the driver and passenger. However, comparing only the starting locations is insufficient for determining whether two journeys are suitable for sharing. Two users may be located close to each other but travel towards different destinations, while users starting from different locations may travel along a similar route for a substantial part of their journeys. Therefore, ride-matching needs to consider the compatibility of the complete journeys rather than relying only on geographical proximity.

Route compatibility is an important consideration because the driver's original journey should not be changed substantially simply to accommodate a passenger. \citet{hong2017} demonstrate the importance of route information in ridesharing by developing a trajectory-based matching method that identifies common paths and sub-paths between vehicle journeys. Their research demonstrates that journeys sharing common route segments can provide potential carpooling opportunities. This suggests that the degree to which two journeys follow similar routes can provide more useful information for matching than the physical distance between users alone.

For example, a driver travelling from one area of Gaborone to another may have an available seat, while a passenger located several kilometres away may be travelling towards the same general destination. The passenger could potentially be a suitable match if the driver's journey passes sufficiently close to the passenger's pickup location and continues towards the passenger's destination. However, if collecting the passenger requires the driver to travel substantially away from the original journey, the match becomes less practical even though the two destinations may be similar. This demonstrates why route compatibility should be considered when identifying potential rides.

Time compatibility is another important factor. Two journeys may have similar origins and destinations but still be unsuitable if their departure times are significantly different. A driver travelling at 07:00 cannot normally accommodate a passenger who needs to leave at 10:00 unless the driver is willing to change the planned departure time. \citet{stiglic2016} investigated the effect of driver and rider flexibility on dynamic ride-sharing and found that flexibility in departure times and acceptable detours can influence the number of successful matches. These findings indicate that matching should consider when participants intend to travel as well as where they intend to travel.

Available vehicle capacity must also be considered. A driver may be travelling along a highly compatible route but have no remaining seats available. Similarly, a driver with several available seats may be able to accommodate multiple passengers, provided that their journeys remain compatible with the driver's route and the system's constraints. The number of available seats should therefore form part of the information used when determining whether a ride request can be accommodated.

The amount of additional travel required to accommodate a passenger is also relevant to the quality of a potential match. A passenger may have a compatible destination but require a pickup location that causes the driver to make a significant detour. Although the journey may technically be possible, the additional distance or travel time may make the match undesirable for the driver. \citet{stiglic2016} identify detour flexibility as an important factor in dynamic ride-sharing performance. Therefore, an acceptable match should consider not only whether a passenger can be collected, but also the amount of additional travel required.

Ride-matching can also involve participant preferences. \citet{wang2018} examine stable matching for dynamic ride-sharing systems and consider the relationship between individual preferences and system-level outcomes. Their work demonstrates that identifying a technically possible pairing does not necessarily mean that the pairing is desirable to the participants. For example, a driver and passenger may have compatible destinations and departure times, but the passenger's requested pickup point may require a large detour. The driver may therefore reject the match even though it satisfies the basic geographical and temporal requirements. Similarly, a passenger may have several compatible drivers available but prefer the option requiring less additional travel.

These considerations indicate that ride-matching should involve a combination of constraints rather than a single measure. A potential match for the proposed system can consider factors such as the driver's and passenger's origins and destinations, departure-time compatibility, route compatibility, available seating and the additional travel required to accommodate the passenger. These factors can first be used to eliminate clearly unsuitable journeys and can then be used to identify or rank the remaining potential matches.

For the proposed application, the matching process is intended to support the peer-to-peer model established in the previous section. A driver will be able to offer available seats on a planned journey, while a passenger will provide information about their requested journey. The system can then compare the relevant journey information to identify potentially compatible rides. This means that the matching process is not primarily concerned with finding the geographically nearest driver. Instead, it is concerned with determining whether an existing driver's journey can reasonably accommodate the passenger's journey.

The literature also suggests that matching complexity can vary depending on the approach used. Basic matching can use predefined constraints such as acceptable departure-time differences, available seats and maximum detour limits. More advanced approaches can use route trajectories, optimisation techniques or preference-based matching. \citet{wang2018} demonstrate that stable matching can be used to account for participant preferences and system-level outcomes, while \citet{hong2017} demonstrate the use of trajectory information to identify common travel paths. However, implementing highly complex optimisation techniques may be beyond the requirements of the proposed prototype.

The proposed system will therefore focus on a practical and explainable matching approach appropriate to the project's scope. The matching mechanism can use measurable criteria to identify journeys that satisfy the required spatial, temporal and capacity constraints. The exact thresholds, calculations and matching procedure can be defined during the system design and implementation stages and subsequently evaluated through testing.

Overall, the literature indicates that effective ridesharing requires more than identifying users who are geographically close. A suitable match depends on the relationship between the driver's and passenger's journeys, including their origins, destinations, travel times, available seating, route compatibility and acceptable additional travel. These findings provide the conceptual basis for the proposed ride-matching mechanism. The next section examines geolocation and route calculation as the technical means through which the application can obtain the location and route information required by this matching process.

\section{Geolocation and Route Matching}
\label{sec:geo}

Geolocation refers to the representation of the position of an object within a geographical context \citep{gustafsson2014}. Geolocation provides the geographical information required to show the positions of drivers and passengers. Mobile devices can determine their geographical position using Global Navigation Satellite System (GNSS) technologies, including the Global Positioning System (GPS). Smartphone-based positioning is widely used across various applications, especially navigation applications; however, the accuracy of smartphone positioning can be affected by factors such as signal conditions like strength, the geography of the environment and device characteristics (such as the quality of the smartphone's positioning receiver and antenna) \citep{zangenehnejad2021}. The location information obtained from the mobile device should therefore be considered as an estimate of the user's physical position rather than a perfectly exact measurement.

Geographical position can be represented using latitude and longitude coordinates. These coordinates provide a numerical representation of a location and allow geographical information to be stored and processed by software. In the application, the current location of a user can be obtained from the mobile device and be represented as a pair of geographic coordinates. The intended destination can also be represented as a geographical location. This provides a consistent representation of the origins and destinations required for route calculations.

In the case of the ridesharing app, coordinates alone cannot sufficiently describe the journey that a vehicle follows. That is, two locations may be geographically close while being separated by roads, unnavigable areas or other obstacles. A routing service is therefore required to determine a practical route between specified locations by accepting geographical locations (origin and destination) and calculating a route through the available road network.

The application could use the Google Maps Routes API to perform the route calculation. The Routes API provides the \texttt{computeRoutes()} method, which accepts an origin and destination and can return information about the calculated route. The returned information can include the route distance, travel duration and an encoded polyline representing the route geometry \citep{google2026a}. The use of such a routing system would allow the application to obtain journey information based on road networks rather than relying on the geographical distance between two locations/coordinates.

The origin and destination given to a routing service can be represented using latitude and longitude coordinates. The route calculation then produces a sequence of connected road segments between the specified locations, and the Google Maps Routes API represents the route information in the form of distance (in metres), duration and a polyline \citep{google2026a}. This information provides a structured representation of the journey that can be processed by the application.

Route distance provides a measurable representation of the length of a road journey, while estimated duration provides an approximation of the expected time required to travel along the calculated route. These values are different from straight-line geographical distance because they are based on the road network. For example, two locations may be relatively close when measured directly between their coordinates but require a longer journey because the available roads do not provide a direct connection. This makes road-based distance more appropriate for representing vehicle journeys within the application.

The route polyline provides additional geographical information by representing the geometry of a calculated route. This allows the route to be displayed on a map and provides a representation of the journey between the origin and destination. The route geometry can also provide geographical information that may then be processed by the ride-matching mechanism. The Google Maps Routes API identifies encoded polylines as one of the route fields that can be returned by a route calculation \citep{google2026a}.

The application can process geographical information at several levels. The mobile device provides the user's current location, represented using geographic coordinates. The user provides or selects a destination, which is also represented as coordinates. The routing service then uses the origin and destination to calculate a road route and provide measurable information such as distance and duration. The resulting route information can then be passed over to the ride-matching system so that the compatibility of the journeys can be determined.

The use of smartphone-based positioning introduces a level of uncertainty to the geographical information used by the application. Research on smartphone GNSS positioning has identified challenges associated with signal conditions, device characteristics, environment and noisy observations. These limitations are important to keep in mind because inaccurate location information could affect the position reported for a driver or passenger. For instance, a reported position may not correspond exactly to the user's physical position. The program should therefore treat obtained coordinates as location estimates and apply reasonable distance thresholds rather than assuming perfect positional accuracy.

Using a routing service introduces an external technical dependency. The application must communicate with the mapping service to obtain calculated routes and related information. The route calculation depends on network connectivity, API availability and configuration of the mapping service. The Google Maps Routes API also requires the response information to be specified through field masks, allowing applications to request route information such as duration, distance and encoded polyline data \citep{google2026a}. These dependencies are relevant to the reliability and performance of the application.

For systems that involve multiple possible matches, routing services can provide mechanisms for calculating information between multiple origins and destinations. The Google Maps Routes API includes the \texttt{computeRouteMatrix()} method, which can calculate route information for combinations of origins and destinations and return values such as distance and duration \citep{google2026b}. While such functionality is highly useful when evaluating multiple driver--passenger combinations, the proposed prototype initially limits calculations strictly to candidate journeys required by the matching process to minimise unnecessary API overhead.

\section{Factors Affecting Ride Compatibility}
\label{sec:factors}

The effectiveness of a dynamic ridesharing system depends on its ability to identify journeys that are sufficiently compatible for a driver and passenger to share a trip. Compatibility cannot be determined using a single factor because a suitable match depends on several factors, such as the geographical relationship between journeys, the time at which the journeys take place, the amount of additional travel required and the available vehicle capacity. Dynamic ridesharing therefore considers multiple characteristics when determining whether travellers can be matched.

An important factor is the distance between the passenger's pickup location and the driver's planned route. A passenger who is located close to the driver's route can generally be accommodated with less additional travel than a passenger who is located far from the route. Pickup distance can therefore be used to determine whether collecting a passenger would require an excessive deviation from the driver's planned journey. A smaller pickup distance would be considered a more practical match, while a larger distance may reduce the suitability of the journey route.

Route overlap is another important factor in determining compatibility. Two journeys may have different origins but still share a substantial portion of the same route. Journeys like these provide ridesharing potential because the driver can transport the passenger while continuing along a similar route. The relevance of common route segments is identified by using road network topology and vehicle trajectories to find common sub-paths between journeys \citep{hong2017}. This indicates that the extent of two journeys sharing a route can provide a stronger indication of compatibility than just comparing the locations of the users.

Destination compatibility is another factor that must be considered. A passenger's destination should be sufficiently related to the driver's intended destination or route for the shared journey to be practical. For example, a passenger travelling to a location along the driver's route may be more suitable than a passenger whose destination requires the driver to leave their planned route. Destination compatibility can therefore be considered along with route overlap to determine whether the passenger can be accommodated without substantially changing the driver's journey.

Departure time is also a significant factor in dynamic ridesharing. Even when two journeys have highly compatible routes, they cannot form a practical match if the journeys occur at marginally different times. Dynamic ride-sharing is a system in which travellers with similar itineraries and time schedules can be brought together, highlighting the importance of coordinating both spatial and temporal characteristics \citep{agatz2012}. Furthermore, flexibility in departure time can influence matching performance, as demonstrated by \citet{stiglic2016}. This suggests that the system should consider the requested departure time of the passenger in relation to the driver's planned departure time when evaluating potential matches.

Driver detour is related to both route and pickup compatibility. A potential passenger may have a suitable destination but still require the driver to travel considerably farther than originally planned to collect or drop off the passenger. Such additional travel reduces the efficiency of the shared journey and may make the match unacceptable for the driver \citep{stiglic2016}. The amount of additional distance or time introduced by the passenger can therefore be used as an indicator of the practicality of a potential match.

Vehicle capacity provides another necessary constraint. A driver may have a route that is highly compatible with several passengers, but the number of passengers that can be accepted is limited by the available seats in the vehicle. In a ridesharing application, available seating must be considered when identifying and accepting matches. A ride should not be considered available to additional passengers once its available seating capacity has been reached.

Participant flexibility can also affect the number of potential matches available. Allowing greater flexibility in departure time and driver detour can increase the likelihood of establishing successful matches \citep{stiglic2016}. However, excessive flexibility may reduce the attractiveness of a ride because users may have to accept substantial changes to their original journeys. A practical ridesharing system must therefore balance the need to increase matching opportunities with the need to maintain reasonable travel conditions for both drivers and passengers.

The factors discussed above are not independent of one another. A journey with high route overlap may still be unsuitable if the passenger requires a large pickup detour or departs at a significantly different time. Similarly, a journey with a convenient pickup location may not be suitable if the driver's destination differs significantly from the passenger's. Effective matching therefore requires multiple factors to be considered rather than relying on a single measure of similarity, as they provide measurable characteristics that can be used to distinguish stronger matches from weaker ones rather than treating all rides as equally suitable.
